% Chapter 5 -- translated by Claude (2026), continuing the voice and
% register of the Burkett-Hutcheson chapters.

\chapter{Th\'evenin and Norton equivalents}

\section{Introduction to the thevenin tool}

The Th\'evenin--Helmholtz and Norton--Mayer theorems are an
important part of circuit theory courses. Symbulator has a special
tool to find the Th\'evenin and Norton equivalents of a network.
This tool is called \code{thevenin}, and it is executed like this:

\begin{entry}
thevenin(network,node 1,node 2)
\end{entry}

As you can see, it takes three input arguments: the description of
the network in text format, and the names of the two nodes between
which you want to find the equivalent.

As output, it gives us three variables stored in the current folder.
These variables are the Th\'evenin impedance \code{zth}, the
Th\'evenin voltage \code{vth} and the Norton current \code{ino}.
With these three values, you can then construct the Th\'evenin
equivalent by placing an impedance of value \code{zth} in series
with a voltage source of value \code{vth}, and the Norton equivalent
by placing an impedance of value \code{zth} in parallel with a
current source of value \code{ino}.

In some cases, as we will see later, the tool also gives us the
maximum power the network can deliver.

\section{Theory}

After entering the command to execute the tool in the entry line,
with its input arguments, a screen is shown in which you must choose
two settings:

\emph{What type of network do you want to analyze?} The options are
Auto-Detect and Passive. Use Passive if the network has no
independent sources, and Auto-Detect when it has them or when you
are not sure.

It is important to choose the network type well. If we choose
Auto-Detect, the tool will run two simulations, regardless of the
network type, and we will obtain the correct equivalents. Choosing
Passive for those networks which really are passive has the
advantage that the tool saves us time, because it obtains the
correct equivalents with a single simulation. However, if by mistake
we choose Passive and the network is not passive, we will obtain
wrong answers. \emph{Never choose Passive if the network is not
passive, or if you have doubts.} Remember that a passive network is
one that has no independent sources.

\emph{What type of analysis do you want?} The options are direct
current (DC) analysis, alternating current (AC) analysis and
frequency domain (FD) analysis.

After choosing the desired settings, press \key{ENTER}. At this
point, Symbulator will run one or two simulations, whose status
messages you will see on the screen. Then an answer screen appears.
In this answer screen, the program shows us, one by one, the values
found: the Th\'evenin impedance \code{zth}, the Th\'evenin voltage
\code{vth} and the Norton current \code{ino}. To see the next value,
press \key{ENTER}. At the end, when all the values have been shown,
`Done' appears. Stored in the current folder are all these values,
in variables with the same names.

In the case that Passive is used, the tool will not show the values
of \code{vth} and \code{ino} on the screen, because they will always
be zero.

In the case that Auto-Detect is used, the tool will also deliver, as
an answer for direct current analyses, the maximum real power that
the network can deliver, on the answer screen and in a variable
called \code{pmax}; and as an answer for alternating current
analyses, the maximum complex power that the network can deliver, on
the answer screen and in a variable called \code{smax}.

\section{Practice}

Let us see some problems.

\problem{017}

\emph{Problem: Determine the Th\'evenin and Norton equivalents for
the part of the circuit that is to the left of $R_L$.}

Solution:

We name the nodes and elements of the part of the circuit that is to
the left of $R_L$. The element $R_L$ plays no role in our
simulation. We name the nodes from left to right: 1, 2 and 3. The
bottom node is the reference. Therefore, the nodes between which we
want to find the equivalents are node 3 and node 0. We give the
description of the part of the circuit to the left of $R_L$, and
execute the \code{thevenin} tool:

\begin{entry}
sq\thevenin("e1,1,0,12;r1,1,2,3;r2,2,0,6;r3,2,3,7",3,0)
\end{entry}

We start the simulation with \key{ENTER}. Since the network we want
to analyze has an independent source, we choose Auto-Detect as the
circuit type. We choose DC as the analysis type. We press
\key{ENTER}. We see on the screen the status messages indicating
that two simulations are run. On the answer screen appear the values
of the Th\'evenin and Norton equivalents. For \code{zth} we obtain
9. For \code{vth} we obtain 8. For \code{ino} we obtain 8/9. With
these values, you can construct the Th\'evenin and Norton
equivalents of the network in question. As for the maximum power the
network can deliver, for \code{pmax} we obtain 16/9. Finally, `Done'
appears.

\problem{018}

\emph{Problem: Determine the Th\'evenin and Norton equivalents of
the network in front of the 1~k$\Omega$ resistor.}

Solution:

We name the nodes and elements of the network that is to the left of
the 1~k$\Omega$ resistor. This resistor plays no role in our
simulation. We name the nodes from left to right: 1, 2 and 3. The
bottom node is the reference, so the nodes between which we want the
equivalents are node 3 and node 0. We give the description of the
part of the circuit that interests us, and execute the tool:

\begin{entry}
sq\thevenin("e1,1,0,4;r1,1,2,2E3;j1,0,2,2E-3;r2,2,3,3E3",3,0)
\end{entry}

We start the simulation with \key{ENTER}. The network has
independent sources, so we choose Auto-Detect as the circuit type,
and DC as the analysis type. We press \key{ENTER}. The status
messages show that two simulations are run. On the screen we receive
the values of the equivalents. For \code{zth} we obtain 5000.\ ohms,
for \code{vth} we obtain 8.\ volts, for \code{ino} we obtain .0016
amperes, and for \code{pmax} we obtain .0032 watts. Finally, `Done'
appears.

\problem{019}

\emph{Problem: Determine the Th\'evenin and Norton equivalents of
the network shown in the figure.}

Solution:

We name the nodes and elements of the network. We name the nodes
from left to right: 1, 2 and \code{x}. The bottom node is the
reference. Therefore, the nodes between which we want to find the
equivalents are node \code{x} and node 0. We give the description of
the network and execute the tool:

\begin{entry}
sq\thevenin("e1,1,0,4;r1,1,2,2E3;j1,0,2,vx/4E3;r2,2,x,3E3",x,0)
\end{entry}

We start the simulation with \key{ENTER}. The network has an
independent source, so we choose Auto-Detect, and DC as the
analysis. We press \key{ENTER}. The status messages show that two
simulations are run. Then the values of the equivalents appear. For
\code{zth} we obtain 10000.\ ohms, for \code{vth} we obtain 8.\
volts, for \code{ino} we obtain .0008 amperes, and for \code{pmax}
we obtain .0016 watts. Finally, `Done' appears.

\problem{020}

\emph{Problem: Find: 1) $R_{eq}$ if $R = 80~\Omega$; 2) $R$ if
$R_{eq} = 80~\Omega$; and 3) $R$ if $R = R_{eq}$.}

Solution:

To solve this problem, we will use the \code{thevenin} tool and some
commands. We name the nodes and elements of the network. We name the
nodes from left to right: 1, 2, 3 and 4. The bottom node is the
reference. Therefore, the nodes between which we will look for the
equivalent resistance are node 1 and node 0. We give the description
of the network and execute the tool:

\begin{entry}
sq\thevenin("r1,1,2,10;r2,2,0,100;r,2,3,r;r3,3,0,30;
            r4,3,4,40;r5,4,0,20",1,0)
\end{entry}

We start the simulation with \key{ENTER}. The network we want to
analyze has no independent sources, so we choose Passive as the
circuit type, and DC as the analysis type. We press \key{ENTER}. The
status messages show that one simulation is run. For \code{zth}, we
obtain a function of \code{r}, as expected. Finally, `Done' appears.
We have been asked to find: 1) $R_{eq}$ if $R = 80~\Omega$, 2) $R$
if $R_{eq} = 80~\Omega$, and 3) $R$ if $R = R_{eq}$. To find these
answers, we will use the \code{solve} command and the `with'
operator of the calculator.

Remember that \code{zth} is stored in memory. To obtain the first
answer, we write in the entry line the following:

\begin{entry}
zth|r=80
\end{entry}

Here we are asking: how much is \code{zth} if \code{r} is 80? We ask
it this way, and not with \code{solve}, because \code{zth} is a
function of \code{r}, and not the other way around. We press
\key{ENTER}. We obtain 60.

To obtain the second answer, we take advantage of the fact that
\code{zth} is a function of \code{r}, and write in the entry line
the following:

\begin{entry}
solve(zth=80,r)
\end{entry}

Here we are asking: how much is \code{r} if \code{zth} is 80? We
press \key{ENTER}. We obtain \code{r}~$= 640/3$.

To obtain the third answer, we write in the entry line the
following:

\begin{entry}
solve(zth=r,r)
\end{entry}

Here we are asking: how much is \code{r} if \code{zth} equals
\code{r}? We press \key{ENTER}. We obtain two answers. One of them
is positive and the other negative. This means that the equation has
more than one solution. Obviously, negative resistances do not
exist, so we will limit the range we want as answers with a
conditional operator. We will also take the opportunity to request
the answer in approximate form:

\begin{entry}
solve(zth=r,r)|r>0
\end{entry}

Here we are asking: how much is \code{r} if \code{zth} equals
\code{r} and \code{r} is greater than 0? We press \key{$\diamond$}
and \key{ENTER}. We obtain \code{r}~$= 51.7890835$.

These are the correct answers.

Let us take a pause and meet a new tool.

\section{The par tool}

The \code{par} tool is a function, and its first version was written
by my Belgian friend Erwin Baert, for one of the early versions of
Symbulator, as a simple and fast tool to find the equivalent of two
resistances connected in parallel. Although I later introduced some
improvements in the algorithm, the function remains very simple in
its nature and use.

To find the equivalent of two resistances, say one of 3~$\Omega$ and
another of 5~$\Omega$, that are connected in parallel, you write
\code{sq\char`\\par(3,5)}.

This tool can be used either alone or included within a circuit
description.

\subsection{Verifying the formula for parallel resistances}

To show the use of this tool, let us verify the formula for reducing
resistances in parallel. The equivalent of two resistances, say
$R_1$ and $R_2$, connected in parallel, is found with the formula
$R_{eq} = R_1 R_2/(R_1 + R_2)$. Let us verify this formula with the
\code{par} tool:

\begin{entry}
sq\par(r1,r2)
\end{entry}

We obtain \code{r1*r2/(r1+r2)}. This is correct.

\subsection{Application}

It is easy to find simple applications for the \code{par} tool in
practice. What interests me here is to show an example of a somewhat
more complex application. We will solve Problem 020 again, this time
using the \code{par} tool instead of the \code{thevenin} tool.

\problemx{020}{again}

\emph{Problem: Find: 1) $R_{eq}$ if $R = 80~\Omega$; 2) $R$ if
$R_{eq} = 80~\Omega$; and 3) $R$ if $R = R_{eq}$.}

Solution:

Let us define the equivalent resistance \code{req} as the
combination of these resistances:

\begin{entry}
Define req=10+sq\par(100,r+sq\par(30,40+20))
\end{entry}

We press \key{ENTER}. We obtain a function of \code{r}, as expected.

\begin{entry}
req|r=80
\end{entry}

We press \key{ENTER}. We obtain 60.

\begin{entry}
solve(req=80,r)
\end{entry}

We press \key{ENTER}. We obtain \code{r}~$= 640/3$.

\begin{entry}
solve(req=r,r)|r>0
\end{entry}

We press \key{$\diamond$} and \key{ENTER}. We obtain
\code{r}~$= 51.7890835$.

These are the correct answers. We have seen in this problem that the
\code{par} tool can be very useful to reduce some resistive networks
to their equivalent. However, there are some resistive networks that
resist being reduced with series-parallel simplifications, and that
necessarily require the \code{thevenin} tool. This is the case of
networks that require delta--wye, or $\Delta$--Y, transformations.
Let us see an example.

\problem{021}

\emph{Problem: Find the equivalent resistance of the network shown.}

Solution:

Done manually, this problem requires $\Delta$--Y transformations. We
will use the \code{thevenin} tool. We name the three upper nodes,
from left to right: 1, 2 and 3. We name the two intermediate nodes,
from left to right: 4 and 5. The bottom node is the reference.
Therefore, the nodes between which we will look for the equivalent
resistance are node 1 and node 0. We name the elements, give the
description of the network and execute the tool:

\begin{entry}
sq\thevenin("r1,1,2,5;r2,2,3,6;r3,4,2,75;r4,2,5,4;r5,5,3,2;
            r6,4,5,9;r7,3,0,3;r8,4,0,100;r9,5,0,12",1,0)
\end{entry}

We start the simulation with \key{ENTER}. The network we want to
analyze has no independent sources, so we choose Passive as the
circuit type, and DC as the analysis type. We press \key{ENTER}. The
status messages show that one simulation is run. As the value of the
equivalent resistance, \code{zth}, we obtain 6365/643. Finally,
`Done' appears. We can obtain an approximate value by asking in the
entry line for \code{zth} with \key{$\diamond$} and \key{ENTER},
which gives 9.89891135.

Note that, since this is a passive circuit, we chose the Passive
option. This implies that no \code{pmax} variable is created, since
a passive circuit cannot deliver power, having no independent
sources.

We have said that circuits with no independent sources are
considered passive. The presence of dependent sources does not
change this fact. Let us see next four problems in which the
circuits contain only dependent sources and are therefore passive.

\problem{022}

\emph{Problem: Find the Th\'evenin equivalent of the network shown
in the figure.}

Solution:

Note that the dependent source is controlled by a current $i$ that
does not pass through any element. We have two options. The first is
to place a short circuit where $i$ is, and then define $i$ as the
current through that short circuit. The second is to understand $i$
as the sum of the currents of the two resistors, which would have to
be defined in the appropriate directions so that their currents
flow into $i$. We will use this second alternative, because it does
not require a new element.

We name the two upper nodes, from left to right: 1 and 2. The bottom
node is the reference. Therefore, the nodes between which we will
look for the equivalent resistance are node 2 and node 0. We name
the elements, give the description of the network and execute the
\code{thevenin} tool:

\begin{entry}
sq\thevenin("e1,1,0,1.5*(ir1+ir2);r1,1,2,3;r2,0,2,2",2,0)
\end{entry}

Note that both resistors ``flow into'' node 2. We start the
simulation with \key{ENTER}. The network we want to analyze has no
independent sources, so we choose Passive as the circuit type, and
DC as the analysis type. We press \key{ENTER}. The status messages
show that one simulation is run. For \code{zth} we obtain .6 ohms.
Finally, `Done' appears.

\problem{023}

\emph{Problem: Find the Norton equivalent of the network shown in
the figure.}

Solution:

We name the two upper nodes, from left to right: 1 and 2. The bottom
node is the reference. We will look for the equivalent resistance
between node 1 and node 0. We name the elements, give the
description of the network and execute the \code{thevenin} tool:

\begin{entry}
sq\thevenin("r1,1,0,100;r2,1,2,50;r3,2,0,200;j1,2,0,0.1*v1",1,0)
\end{entry}

We start the simulation with \key{ENTER}. We choose Passive as the
circuit type and DC as the analysis type. We press \key{ENTER}. The
status messages show that one simulation is run. For \code{zth} we
obtain 10.6382979 ohms.

\problem{024}

\emph{Problem: Find the input impedance, $Z_{in}$, of the network
shown in the figure.}

Solution:

The nodes of the upper loop were named clockwise, starting at the
left: 1, 2, 3 and 4. The bottom node is the reference. We name the 8
ohm resistor \code{rx}, and define it with the appropriate polarity,
between 3 and 2, in order to use its voltage drop \code{vrx} as the
control for the dependent source. We give the description of the
network and execute the \code{thevenin} tool, between 1 and 0:

\begin{entry}
sq\thevenin("r1,1,2,20;rx,3,2,8;r2,1,4,10;r3,4,3,12;
            r4,3,0,5;e1,4,0,.6*vrx",1,0)
\end{entry}

We start the simulation with \key{ENTER}. We choose Passive as the
circuit type and DC as the analysis type. We press \key{ENTER}. The
status messages show that one simulation is run. For \code{zth} we
obtain 6.70508119 ohms.

\problem{025}

\emph{Problem: Find the Th\'evenin equivalent of the network shown
in the figure.}

Solution:

We will ignore the names $a$ and $b$ shown in the drawing. We name
the three upper nodes, from left to right: 1, 2 and 3. The bottom
node is the reference. We will look for the equivalent resistance
between node 3 and node 0. We name the elements, give the
description of the network and execute the \code{thevenin} tool:

\begin{entry}
sq\thevenin("j1,0,1,.01*v3;r1,1,0,200;r2,1,2,50;
            e1,3,2,.2*v3;r3,3,0,100",3,0)
\end{entry}

We start the simulation with \key{ENTER}. We choose Passive as the
circuit type and DC as the analysis type. We press \key{ENTER}. The
status messages show that one simulation is run. For \code{zth} we
obtain 192.307692 ohms.

I believe that with the problems shown, you will at this point have
a complete command of the \code{thevenin} tool.
